\documentclass[a4paper,11pt]{article}

\usepackage{fontspec}
\setmainfont{Carlito}
\usepackage[margin=0.6in, top=0.5in, bottom=0.5in]{geometry}
\usepackage{titlesec}
\usepackage{enumitem}
\usepackage{hyperref}
\usepackage{xcolor}
\usepackage{microtype}
\usepackage{needspace}

\hyphenpenalty=10000 \exhyphenpenalty=10000 \emergencystretch=3em

\definecolor{headernavy}{RGB}{12, 83, 166}
\definecolor{linkblue}{RGB}{0, 102, 204}
\definecolor{muted}{RGB}{96, 102, 110}

\hypersetup{colorlinks=true, linkcolor=linkblue, urlcolor=linkblue}

\titleformat{\section}
  {\color{headernavy}\normalsize\bfseries}
  {}{0em}{\MakeUppercase}[\color{headernavy}\titlerule]

\titlespacing{\section}{0pt}{8pt}{4pt}

\setlist[itemize]{noitemsep, topsep=1pt, leftmargin=1.3em, partopsep=0pt, parsep=1pt}

\pagestyle{empty}
\setlength{\parindent}{0pt}
\linespread{1.02}

% ---- building blocks: bold = organisation/project, italic = role/stack, grey = dates/places ----
\newcommand{\role}[4]{% org, location, dates, title
  \textbf{#1}{\color{muted}\,\textperiodcentered\,\textit{#2}}\hfill{\color{muted}#3}\\
  \textit{#4}\par}
\newcommand{\roleNoLoc}[3]{\textbf{#1}\hfill{\color{muted}#2}\\ \textit{#3}\par}
\newcommand{\project}[3]{% name, stack, right-hand tag
  \textbf{#1}\hfill#3\\ {\color{muted}\small\textit{#2}}\par}
\newcommand{\gh}[1]{\href{#1}{\small GitHub}}
\newcommand{\badge}[1]{{\color{headernavy}\small\textbf{#1}}}
\newcommand{\skillgroup}[2]{{\color{headernavy}\textbf{#1}}\\ #2\par\vspace{3pt}}
\newcommand{\gap}{\vspace{5pt}}
\newcommand{\num}[1]{\textbf{#1}}   % result numbers only


\begin{document}

% ================= HEADER =================
\begin{center}
  {\LARGE\bfseries Thiruvel Andagurunathan Pandian} \\[3pt]
  {\color{headernavy}\small\textbf{Data Scientist}} \\[3pt]
  {\small
  Bristol, UK\nobreak\ \,\textbar\, \href{mailto:apthiruvel@gmail.com}{apthiruvel@gmail.com}\nobreak\ \,\textbar\, \href{tel:+447741016954}{+44 7741 016954}\nobreak\ \,\textbar\, \href{https://www.linkedin.com/in/thiruvel-a-p/}{https://www.linkedin.com/in/thiruvel-a-p/}\nobreak\ \,\textbar\, \href{https://github.com/Thiruvel-AP}{https://github.com/Thiruvel-AP}\nobreak\ \,\textbar\, \href{https://thiruvel-ap.github.io/PersonalPortfolio/}{https://thiruvel-ap.github.io/PersonalPortfolio/}}\par\vspace{4pt}
  \noindent{} Right to Work in UK \,\textbar\, Available to join Immediately
\end{center}

% ================= PROFILE =================
\needspace{5\baselineskip}\section{Profile}
Builds the modelling a brief needs and reports its limits plainly: gradient boosting (XGBoost churn), LSTM time series forecasting with time ordered validation, text clustering (SBERT, K-Means, DBSCAN) and 11,459 per gene models (train Pearson 0.807 vs test 0.545). Works with clients: scoped a loosely defined brief at 180 Degrees Consulting into a stakeholder demo (five star rating) and contributed to the client\textquotesingle{}s production MCP server, now in live use. Deploys on AWS and GCP (Lambda, Cloud Run) with CI/CD and MLflow experiment tracking. Builds deployable ML models: forecasting (an LSTM time series service), propensity modelling (an XGBoost churn classifier with Bayesian tuning and SMOTE) and customer segmentation (k-means, HAC and DBSCAN).

% ================= SKILLS =================
\needspace{5\baselineskip}\section{Skills}
\skillgroup{Machine Learning, Deep Learning \& Statistics}{Predictive modelling (gradient boosting with XGBoost, random forest, scikit-learn), clustering (K-Means, HAC, DBSCAN, Leiden), deep learning (PyTorch, TensorFlow/Keras, LSTM, GRU, BiRNN, CNN), time series forecasting (cyclic time encoding, TimeSeriesSplit), anomaly detection (Isolation Forest), class imbalance handling (SMOTE), hyper parameter tuning (grid search, two stage Bayesian), model evaluation (held out, cross validation, Pearson, MCC), statistical inference and experimental design \textit{(confidence intervals, bootstrap, hypothesis testing; MSc coursework)}, MLflow experiment tracking}
\skillgroup{Python, Data Quality \& Cloud}{Python, pandas, NumPy, SciPy, Matplotlib, Seaborn, Plotly, SQL, PostgreSQL, R (tidymodels), identifier crosswalks, parquet, reproducible pipelines, AWS, GCP, Docker, Kubernetes (AWS EKS), serverless deployment (Lambda, Cloud Run), containerisation, GitHub Actions CI/CD, FastAPI, TensorFlow Serving, pytest, Git}
\skillgroup{NLP, LLMs, Embeddings \& Evaluation}{NLP (text preprocessing, TF-IDF, SBERT embeddings, text clustering), Hugging Face Transformers, large language models (Claude, OpenAI, open weight via Ollama), prompt engineering, RAG with embeddings (ChromaDB), agents (LangChain, LangGraph), team LoRA fine tune of Llama, LLM evaluation (hallucination rate analysis, deterministic guardrails)}
\skillgroup{Client Delivery \& Communication}{Scoping ambiguous problems with clients, stakeholder demos, working with client engineering teams, code review, transparent reporting of model limitations, root cause analysis, technical documentation}

% ================= EXPERIENCE =================
\needspace{5\baselineskip}\section{Work Experience}
\needspace{4\baselineskip}
\role{AstraZeneca}{Bristol, UK (Remote)}{May 2026 – Sep 2026}{Data Scientist}
\begin{itemize}
  \item Needed an LLM a research team could trust before it shipped: designed a deterministic hallucination guard architecture and fail safe data resolution via live Ensembl/HGNC APIs with a local \num{19,213 gene parquet fallback} and circuit breakers.
  \item Wrote the dissertation\textquotesingle{}s AI agent sections covering architecture, implementation and evaluation, including a documented hallucination rate analysis and test case outcomes.
  \item Diagnosed where it struggled with a three tier, \textasciitilde{}65 test evaluation harness (tool level → API → live) and hallucination rate analysis, debugging failures layer by layer, and used the same harness to validate the Groq to Bedrock migration after a rate limit evaluation.
\end{itemize}
\gap
\needspace{4\baselineskip}
\role{180 Degrees Consulting Bristol}{Bristol, UK (Remote)}{Sep 2025 – Dec 2025}{Data Consultant (AI Engineering \& AI System)}
\begin{itemize}
  \item Faced a loosely defined problem from a Madrid based startup client: proposed and prototyped an MCP based agentic tool in Python (FastMCP, Ollama gpt-oss:20b, PostgreSQL, GitHub Actions CI), presented in a live stakeholder demo.
  \item Translated ambiguous business requirements into a production ready technical blueprint and owned the solution within a cross functional consulting team, from architecture design through stakeholder presentation.
  \item Proposed the request flow: a user\textquotesingle{}s text goes through the API and a customised prompt to the MCP server, where one LLM instance writes GraphQL queries against the client\textquotesingle{}s Convex database and a second instance runs the statistical analysis before the result returns to the user.
\end{itemize}
\gap
\needspace{4\baselineskip}
\role{Freelance}{Chennai, India}{Jan 2025 – Aug 2025}{AI/ML Engineer}
\begin{itemize}
  \item Designed reproducible scikit-learn pipelines with feature engineering (temporal features, interaction terms, target encoding), model selection, cross validation and hyperparameter tuning across classification and regression tasks.
  \item Working independently for India based startups and individual clients, built agentic RAG workflows (LangChain, LangGraph, ChromaDB) on GCP Cloud Run, NLP pipelines (Hugging Face Transformers, Ollama) and scikit-learn pipelines, served ML models as REST APIs (FastAPI, Docker), and contributed to a team LoRA fine tune of a Llama model.
\end{itemize}
\gap
\needspace{4\baselineskip}
\role{Avasoft}{Chennai, India}{Jan 2023 – Sep 2024}{Software Development Engineer (Jun 2023 – Sep 2024); Software Development Engineer Trainee (Jan 2023 – May 2023)}
\begin{itemize}
  \item Integrated Twelve Data (real time market feeds) and Plaid (bank account linkage, transaction imports) APIs into a Flutter FinTech mobile app.
  \item Secured backend communication with SSL pinning and integrated AWS Cognito for user authentication with single sign on in the FinTech app.
  \item Built ecommerce and ticketing UIs in Swift/SwiftUI with product search, cart logic, order tracking and booking management.
\end{itemize}

% ================= PROJECTS =================
\needspace{5\baselineskip}\section{Projects}
\needspace{4\baselineskip}
\setbox0=\hbox{\textbf{llama3 8b qlora math t4: QLoRA fine tuning of Llama 3.1 8B for maths under a token budget}\hspace{1em}\url{https://github.com/Thiruvel-AP/llama3-8b-qlora-math-t4}}%
\setbox2=\hbox{{\small\textit{Python, Unsloth, PEFT, TRL, Hugging Face, PyTorch, sentence transformers, scikit-learn}}\hspace{1em}{\small\url{https://github.com/Thiruvel-AP/llama3-8b-qlora-math-t4}}}%
\ifdim\wd0>\linewidth
\ifdim\wd2>\linewidth
\textbf{llama3 8b qlora math t4: QLoRA fine tuning of Llama 3.1 8B for maths under a token budget}\\ {\color{muted}\small\textit{Python, Unsloth, PEFT, TRL, Hugging Face, PyTorch, sentence transformers, scikit-learn}}\\ \hspace*{\fill}{\small\url{https://github.com/Thiruvel-AP/llama3-8b-qlora-math-t4}}\par
\else
\textbf{llama3 8b qlora math t4: QLoRA fine tuning of Llama 3.1 8B for maths under a token budget}\\ {\color{muted}\small\textit{Python, Unsloth, PEFT, TRL, Hugging Face, PyTorch, sentence transformers, scikit-learn}}\hfill{\small\url{https://github.com/Thiruvel-AP/llama3-8b-qlora-math-t4}}\par
\fi
\else
\project{llama3 8b qlora math t4: QLoRA fine tuning of Llama 3.1 8B for maths under a token budget}{Python, Unsloth, PEFT, TRL, Hugging Face, PyTorch, sentence transformers, scikit-learn}{\url{https://github.com/Thiruvel-AP/llama3-8b-qlora-math-t4}}
\fi
\begin{itemize}
  \item Needed to know what fine tuning can change in an 8B model on a free 15 GB GPU: reported the limits openly: \num{accuracy on completed answers barely moved (50\% vs 52\%)}, so the gain came from finishing answers inside the token budget rather than stronger reasoning.
  \item Traced out of memory errors to several models loaded in one runtime, not the KV cache.
\end{itemize}
\gap
\needspace{4\baselineskip}
\setbox0=\hbox{\textbf{Bitcoin Forecasting: LSTM time series service}\hspace{1em}\url{https://github.com/Thiruvel-AP/BitcoinForecasting}}%
\setbox2=\hbox{{\small\textit{Python, TensorFlow, TF Serving}}\hspace{1em}{\small\url{https://github.com/Thiruvel-AP/BitcoinForecasting}}}%
\ifdim\wd0>\linewidth
\ifdim\wd2>\linewidth
\textbf{Bitcoin Forecasting: LSTM time series service}\\ {\color{muted}\small\textit{Python, TensorFlow, TF Serving}}\\ \hspace*{\fill}{\small\url{https://github.com/Thiruvel-AP/BitcoinForecasting}}\par
\else
\textbf{Bitcoin Forecasting: LSTM time series service}\\ {\color{muted}\small\textit{Python, TensorFlow, TF Serving}}\hfill{\small\url{https://github.com/Thiruvel-AP/BitcoinForecasting}}\par
\fi
\else
\project{Bitcoin Forecasting: LSTM time series service}{Python, TensorFlow, TF Serving}{\url{https://github.com/Thiruvel-AP/BitcoinForecasting}}
\fi
\begin{itemize}
  \item Needed forecasts validated on time ordered data: trained a forecasting LSTM with cyclic date encoding and TimeSeriesSplit grid search tuning, then deployed it behind a TensorFlow Serving REST API.
\end{itemize}

% ================= EDUCATION =================
\needspace{5\baselineskip}\section{Education}
\needspace{3\baselineskip}
\roleNoLoc{University of Bristol}{Sep 2025 – Sep 2026}{MSc Data Science, Distinction \textit{(Expected)}}
{\color{muted}\small Large scale Data Engineering, Statistical Computing \& Empirical Methods, AI \& Text Analytics, Visual Analytics, Data Science Methods \& Practices, Technology, Innovation, Business \& Society}\par
\gap
\needspace{3\baselineskip}
\roleNoLoc{St. Joseph\textquotesingle{}s College of Engineering}{Aug 2019 – Apr 2023}{B.E. Electrical and Electronics Engineering, 91\%, First Class with Distinction}

% ================= CERTIFICATIONS & ACHIEVEMENTS =================
\needspace{5\baselineskip}\section{Certifications \& Achievements}
\begin{itemize}
  \item Model Context Protocol: Advanced Topics, Anthropic, \textit{Jun 2026}
  \item Introduction to Model Context Protocol, Anthropic, \textit{Jun 2026}
  \item AWS Academy Cloud Foundations, \textit{Nov 2025}
  \item 180 Degrees Consulting Data Science Consultant Certificate, \textit{Dec 2025}
  \item Bristol PLUS Award, \textit{Mar 2026}
\end{itemize}

\end{document}